% tonks_paper03_v1.tex
% Hard-rod gases resist Burgers shock formation under integrability breaking.
%
% v1 (2026-05-11): negative result, paired with Hübner-Doyon 2024.

\pdfminorversion=7
\documentclass[aps,pre,twocolumn,superscriptaddress,longbibliography,floatfix]{revtex4-2}

\usepackage{amsmath,amssymb,amsfonts,amsthm}
\usepackage{graphicx}
\usepackage{bm}
\usepackage{hyperref}
\usepackage{physics}
\usepackage{mathtools}

\begin{document}

\title{Mass-disordered hard rods resist Burgers shock formation:\\
a microscopic study on Riemannian rings}

\author{Hector Medel}
\email{hjmedel@gmail.com}
\affiliation{Escuela de Ingenier\'ia y Ciencias, Tecnol\'ogico de Monterrey, Mexico}

\date{\today}

\begin{abstract}
The recent rigorous result of H\"ubner and Doyon
[arXiv:2406.18322, 2024] establishes that the integrable
Lieb--Liniger gas does not develop Burgers shocks at any finite
time. We ask whether the canonical microscopic mechanism for
restoring shocks --- breaking the conserved velocity multiset
by binary mass disorder --- suffices in a closely related
classical setting: the one-dimensional hard-rod gas on a closed
Riemannian curve. Event-driven molecular dynamics across four
initial-condition families (sinusoidal density, sinusoidal
velocity, right-moving sound mode, and step velocity) at
amplitudes up to $A=0.8$, step amplitudes up to
$U_0/c_s\approx 1.8$ (well into the supersonic regime), and mass
disparities $m_{\max}\in\{5,10,20,100\}$ spanning the limit of
nearly-immobile impurities show \emph{no} Burgers-shock
cascade in any configuration, despite the integrability-breaking
time scale
$\tau_{\rm break}\sim 0.05$ being two orders of magnitude shorter
than the naive shock time $t^*_{\rm shock}\sim 4$--$7$. The
mode-$n$ density amplitude $|\hat\rho_{k_n}|$ decays faster than
the Burgers $1/n$ signature even for the most singular initial
condition; the closest approach is a transient $n=2$ enhancement
generated by second-order nonlinear coupling between
linearly-populated odd modes, with no extension to
higher-order cascade. Arc-length universality
established for the linear regime in the companion
paper~\cite{PaperI} persists in the nonlinear regime: data
on a circular ring and on a wavy ring with $44\%$ larger
perimeter are statistically indistinguishable. A phase-space
moment analysis with bin convergence, chunk-bootstrap errors,
and kurtosis check of local thermodynamic equilibrium confirms
that the second-moment closure fails for sharp velocity
gradients --- under Euler with both global-$\eta$ and local
Tonks pressures and under Navier--Stokes with any uniform
viscosity; total energy is conserved exactly while coherent
kinetic energy decays on a time scale of order unity into
thermal motion, demonstrating that the no-shock result is not
attributable to absent dissipation. We interpret the results in operational rather than
nomenclatural terms: on the time scales and amplitudes
accessible, the mass-disordered hard-rod gas does not develop
the Fourier signature of a Burgers shock, and no second-moment
hydrodynamic closure --- Euler, Tonks-corrected, or
Navier--Stokes with uniform viscosity --- reproduces the
observed dynamics for sharp gradients. The natural candidate
for an effective theory remains the generalised hydrodynamic
equation with rapidity-resolved quasiparticle densities, but
we make no positive claim of GHD-protection beyond consistency.
\end{abstract}

\maketitle

\section{Introduction}

The hard-rod gas in one dimension occupies a privileged position
in non-equilibrium statistical mechanics: equilibrium structure
is classical \cite{Tonks1936,SalsburgZwanzig1953,LiebMattis1966};
ergodicity in the infinite-volume limit was proved by Aizenman,
Goldstein and Lebowitz \cite{Aizenman1975}; the dynamic structure
factor was computed by Lebowitz, Percus and Sykes
\cite{Lebowitz1968}; and the hydrodynamic limit was placed on
rigorous footing by Boldrighini, Dobrushin and Sukhov
\cite{BDS1983}. The modern Generalised Hydrodynamics (GHD)
framework \cite{DoyonSpohn2017,Doyon2017Lectures,DenardisDoyon2018,
Bulchandani2018,GHDperspective2025} subsumes both, identifying
$v_{\rm eff}=v/(1-\eta)$ as the dressed velocity that controls
all linear-response observables. Experimentally, the hard-rod
description has direct relevance to single-file colloidal channels
\cite{Wei2000,LinCui2015}, ultracold-atom realisations of the
Lieb--Liniger gas where GHD has been quantitatively tested
\cite{Schemmer2019,Malvania2021,Cataldini2022}, and to
classical analogues such as ballistic granular chains. The
question of how robustly the integrable phenomenology survives
controlled perturbations is therefore of operational as well as
theoretical interest.

A central feature of these integrable one-dimensional gases is
the \emph{absence of Burgers shocks}. The classical sound mode
of an interacting fluid would develop a shock at finite time
under finite-amplitude perturbation through Burgers nonlinearity.
This does not happen in integrable systems: the conserved
velocity (or rapidity) multiset prevents the cascade of
energy to high wavenumbers that produces a shock. The recent
rigorous result of H\"ubner and Doyon \cite{HubnerDoyon2024}
proves this for the Lieb--Liniger model by constructing a new
quadrature for the GHD equation and showing that no shock
develops at any finite time.

The natural question is: \emph{what microscopic ingredient is
needed to restore Burgers shocks?} The canonical answer in
statistical physics is to break the integrability. For a
hard-rod gas, the simplest controlled breaking is binary mass
disorder: a fraction of rods has a different mass, so elastic
collisions no longer simply exchange velocities and the velocity
multiset is no longer conserved. In a companion paper
\cite{PaperII}, we have shown that this mass disorder
\emph{does} restore the sound-wave content of the linear
response: the integrable Tonks gas exhibits no sound waves
in the density-mode autocorrelation, while a mass disparity as
small as $20\%$ produces resolved underdamped sound waves with
a frequency parametrically distinct from the standard
hydrodynamic prediction. It is therefore natural to expect that
the same mass disorder, applied at finite amplitude, would
produce Burgers shocks. The expectation rests on the standard
hydrodynamic logic: linear-regime sound waves plus nonlinear
self-steepening yield shocks in any second-moment closed
fluid. We will show that this naive promotion of the linear
result fails: sound-wave restoration in linear response and
shock absence in nonlinear response are simultaneously
realized, signalling that second-moment closure itself
breaks down.

Here we show that it does not. Across a scan of four
initial-condition families, amplitudes up to $A=0.8$, step
amplitudes spanning the sub- and supersonic regimes
($U_0/c_s$ up to $1.8$), mass disparities
$m_{\max}\in\{5,10,20,100\}$,
and Riemannian-curve geometries, the mass-disordered hard-rod
gas does not develop the Fourier signature of a Burgers shock --- this despite the
fact that the integrability-breaking time scale
$\tau_{\rm break}\sim 0.05$ is two orders of magnitude shorter
than the naive shock time $t^*_{\rm shock}\sim 4$--$7$
(Sec.~\ref{sec:integrability}, Sec.~\ref{sec:timescales}). The
closest approach to shock-like behaviour is a transient $n=2$
enhancement in step-IC runs, traceable to second-order
nonlinear coupling between linearly-populated odd modes, with
no extension to a higher-order cascade. A phase-space moment analysis with bin convergence
and chunk-bootstrap errors demonstrates that the second-moment
closure (Euler with global-$\eta$ or local Tonks pressure, and
Navier--Stokes with any uniform viscosity) fails for sharp
velocity gradients: the residual exceeds the continuity floor by
an order of magnitude at the working resolution and grows with
refinement.

The contributions are:
\emph{(i)} a microscopic verification of the H\"ubner--Doyon
result in pure hard rods at finite amplitude
(Sec.~\ref{sec:tonks});
\emph{(ii)} a negative result that mass disorder alone does not
restore Burgers shocks within the accessible time, amplitude,
and mass-disparity windows, with the second-order origin of the
$n=2$ enhancement isolated from a Burgers-cascade interpretation
(Sec.~\ref{sec:massdisorder}--\ref{sec:step});
\emph{(iii)} a nonlinear-regime confirmation of the arc-length
universality theorem of Ref.~\cite{PaperI}
(Sec.~\ref{sec:curved});
\emph{(iv)} a phase-space moment diagnostic, combining bin
convergence, chunk-bootstrap errors, kurtosis test of local
thermodynamic equilibrium, viscous-Navier--Stokes closure sweep,
and per-species energy partition, showing that no second-moment
hydrodynamic closure reproduces the observed dynamics for sharp
gradients while a finite coarse-grained dissipation rate is
nevertheless present (Sec.~\ref{sec:phase}); and
\emph{(v)} the empirical observation that the density-mode
spectrum is statistically indistinguishable across a 20-fold
range of mass disparity ($m_{\max}=5,20,100$), an open
constraint on any candidate effective theory
(Sec.~\ref{sec:heavy}, Sec.~\ref{sec:discussion}).

\section{Setup}\label{sec:setup}

We consider $N$ hard rods of diameter $\sigma$ on a smooth
closed curve $\gamma:S^1\to\mathbb R^d$ with arc-length element
$ds=\sqrt{g(\varphi)}\,d\varphi$ and total length
$L=\oint\sqrt{g}\,d\varphi$. The packing fraction is
$\eta=N\sigma/L\in[0,1)$, and we work at temperature $k_BT$ with
thermal velocity $v_{\rm th}=\sqrt{k_BT/m}$ ($m=k_BT=1$). The
configurational integral and arc-length reduction are described
in Refs.~\cite{PaperI,PaperII}. In arc-length coordinates the
dynamics is the flat Tonks gas; geometry enters only through
$L$ at the linear-response level.

\subsection{Initial conditions}

We test four IC families, all of the form
\begin{equation}
\rho(s,0)=\rho_0[1+A_\rho \cos(k s)],\quad
v_i(0)=\xi_i^{\rm thermal}+u_{\rm coh}(s_i),
\end{equation}
where $k=2\pi/L$ is the fundamental wavenumber,
$\xi_i^{\rm thermal}\sim\mathcal{N}(0,k_BT/m_i)$, and the
coherent velocity field $u_{\rm coh}$ takes the following forms:

\noindent
\textbullet\ Density-only ($\alpha$): $u_{\rm coh}(s)=0$,
$A_\rho>0$.\\
\textbullet\ Velocity wave ($\beta$): $A_\rho=0$,
$u_{\rm coh}(s)=u_0\cos(ks)$.\\
\textbullet\ Sound mode ($\gamma$): $A_\rho>0$,
$u_{\rm coh}(s)=c_s\,A_\rho\cos(ks)$, with
$c_s=v_{\rm th}/(1-\eta)$.\\
\textbullet\ Step velocity ($\delta$, canonical Burgers Riemann
setup): $A_\rho=0$,
$u_{\rm coh}(s)=u_0\,\mathrm{sgn}[\cos(ks)]$.

The constraint $\sum_i m_i v_i=0$ is enforced for all setups.
Positions are constructed by Newton inversion of the cumulative
density. For mass-disordered runs, half the rods have $m=1$,
half $m=m_{\max}$, randomly shuffled per realisation.

\subsection{Observables}

The mode-$n$ density Fourier component is
\begin{equation}
\hat\rho_{k_n}(t) = \frac{1}{N}\sum_{i=1}^N
e^{-i k_n s_i(t)},
\quad k_n=2\pi n/L,
\end{equation}
ensemble-averaged over independent realisations. We compute
both the cosine and sine components separately and recover the
magnitude $|\hat\rho_{k_n}|=\sqrt{\langle\cos\rangle^2+
\langle\sin\rangle^2}$, an unbiased estimator (cf.\
App.~\ref{app:estimator}). For phase-space moments, we
accumulate spatially-binned counts, mean velocity and kinetic
temperature:
$\rho(s,t)$, $u(s,t)$, $T(s,t)$ on $80$ uniform bins in $s$.

\subsection{Numerical methods}

Event-driven simulations evolve $\{(s_i,v_i)\}_{i=1}^N$ exactly
between collisions; pair velocities are exchanged in the
equal-mass case and updated by the elastic momentum-conserving
collision rule
$v_i'=[(m_i-m_j)v_i + 2 m_j v_j]/(m_i+m_j)$ in the binary
mixture. Eleven datasets are used:
\textsc{tonksA} (pure Tonks finite-$A$, $A\in\{0.1,0.3,0.5\}$,
IC types $\alpha,\gamma$),
\textsc{secIV} (mass-disorder $m_{\max}=5$ at the same
amplitudes and IC types),
\textsc{A08} ($A=0.8$ push at $m_{\max}=5$),
\textsc{stepUB} (step IC, $m_{\max}=5,10$,
$U_0\in\{0.3,0.5\}$, unbiased estimator),
\textsc{curved} (step IC, $m_{\max}=5$, $U_0=0.5$ on wavy
metric $r(\varphi)=R[1+0.3\cos(3\varphi)]$),
\textsc{phasespace} (spatial moments, three IC families
$\alpha,\delta,\gamma$, $n_{\rm bins}=80$),
\textsc{refinement} (spatial moments for step IC at
$n_{\rm bins}\in\{40,80,160\}$, $6000$ seeds in $30$
chunks for chunk-bootstrap),
\textsc{supersonic} (step IC sweep
$U_0\in\{0.5,1.0,1.5,2.0\}$ at $m_{\max}=5$,
$4000$ seeds, density Fourier modes only),
\textsc{kurtosis} (phase-space moments up to $v^4$, step IC
$U_0=0.5$, $m_{\max}=5$, $4000$ seeds, for the local-thermal-%
equilibrium check of Sec.~\ref{sec:phase}),
\textsc{entropy} (per-species spatial moments, step IC
$U_0=0.5$, $m_{\max}=5$, $4000$ seeds, for the energy-partition
diagnostic of Sec.~\ref{sec:phase}),
\textsc{heavymass} (step IC $U_0=0.5$,
$m_{\max}\in\{20,100\}$, $4000$ seeds each).
All runs use $N=400$, $\eta=0.10$, $3000$--$10\,000$ seeds per
configuration.

\subsection{Integrability status of the binary mixture}
\label{sec:integrability}

A central question for the present work is whether the binary
mass-disordered hard-rod gas is integrable. The equal-mass Tonks
gas is integrable in the strongest sense: collisions exchange
velocities, the velocity multiset is exactly conserved, and the
Jepsen mapping \cite{Jepsen1965} reduces the dynamics to a
non-interacting gas. The binary mixture
$m_i\in\{1,m_{\max}\}$ with $m_{\max}\neq 1$ \emph{breaks} this
exact structure: a collision between unequal masses transfers a
fraction $4 m_1 m_2/(m_1+m_2)^2$ of the kinetic energy across
species, and the per-species velocity multisets are no longer
individually conserved. The microcanonical conserved quantities
reduce to the global $\{N_1, N_2, P_{\rm tot}, E_{\rm tot}\}$.

The natural timescale for this integrability-breaking is the
inter-species collision time. For binary fraction $\phi=1/2$,
mean relative speed $\bar v\sim\sqrt{k_B T (m_1+m_2)/(m_1 m_2)}$,
and packing fraction $\eta$,
\begin{equation}
\tau_{\rm break} \sim \frac{(1-\eta)\,L}{N\,\bar v}\cdot
\frac{1}{2\phi(1-\phi)} = \frac{2(1-\eta)L}{N\bar v}.
\label{eq:taubreak}
\end{equation}
For our parameters ($L\approx 12.57$, $N=400$, $\eta=0.1$,
$m_{\max}=5$, $k_B T=1$): $\bar v=\sqrt{6/5}\approx 1.10$
giving $\tau_{\rm break}\approx 5.1\times 10^{-2}$. For
$m_{\max}=10$, $\bar v\approx 1.05$ and
$\tau_{\rm break}\approx 5.4\times 10^{-2}$. In both cases
$\tau_{\rm break}$ is two orders of magnitude shorter than the
naive shock time
$t^*_{\rm shock}\sim 4$--$7$ (Sec.~\ref{sec:timescales}), so
integrability-broken dynamics is established well before any
shock could form. The trivial alternative ``no shocks form
because the system is still effectively integrable on the
simulated timescale'' is ruled out.

We emphasize that not all integrability-breaking channels are
exhausted by binary mass disorder: a soft (finite-range)
inter-particle potential, dissipative or inelastic collisions,
or a multi-mass (rather than binary) distribution would each
provide qualitatively different kernels. Three-body collisions
have measure zero for point-like hard rods and finite collision
duration is excluded by construction; these channels are
specific to soft-potential models. Our results address binary
elastic mass disorder specifically.

\section{Theoretical background}\label{sec:theory}

\subsection{Covariant Burgers}

The covariant inviscid Burgers equation on $(M,g)$ reads
\begin{equation}
\partial_t u + u\,\partial_\varphi u + \Gamma(\varphi)\,u^2 = 0,
\label{eq:burgers}
\end{equation}
with Christoffel symbol $\Gamma=g'/(2g)$ \cite{KhesinMisiolek2007};
its viscous version adds a term $\nu\,\partial_\varphi^2 u$ on
the right-hand side and is the model of interest when a finite
microscopic dissipation is present.
Equation~\eqref{eq:burgers} predicts shock formation at time
$t^* \sim 1/(k\,u_0)$ for a sinusoidal initial velocity profile
of amplitude $u_0$. The characteristic shock signature is a
velocity-mode cascade $|\hat u_{k_n}(t^*)|\sim u_0/n$ at all
$n$. Translated to density modes via continuity
(Sec.~\ref{sec:rhoVSu}),
\begin{equation}
\frac{|\hat\rho_{k_n}(t^*)|}{|\hat\rho_{k_1}(t^*)|}
\;\approx\; \text{const} \sim 1\quad
\text{at } t\!\approx\!t^*
\quad\text{(Burgers shock signature)},
\label{eq:burgersratio}
\end{equation}
i.e.\ a true Burgers shock produces a density cascade that is
flat in $n$. In the alternative weighted form
$n|\hat\rho_{k_n}|/|\hat\rho_{k_1}|$, this is equivalent to a
linearly growing ratio $\sim n$. For comparison, a smooth
(phase-mixed) wave gives $|\hat\rho_{k_n}|/|\hat\rho_{k_1}|$
decaying faster than $1/n$.

\subsection{Density modes as cascade witnesses}
\label{sec:rhoVSu}

Burgers' equation is usually stated for the velocity field; our
data record density Fourier modes
$\hat\rho_{k_n}(t)=N^{-1}\sum_i e^{-ik_n s_i(t)}$. We show here
that a Burgers cascade in $u$ is faithfully reflected in $\hat\rho$.

Continuity gives, after linearization in $\delta\rho=\rho-\rho_0$,
\begin{equation}
\partial_t \hat\rho_{k_n} + i k_n \rho_0 \,\hat u_{k_n}
= -i k_n\,\widehat{(\delta\rho\,u)}_{k_n}.
\label{eq:continuity-fourier}
\end{equation}
For a Burgers shock developing from a sinusoidal or step
velocity perturbation of characteristic amplitude $u_0$, the
velocity cascade obeys $|\hat u_{k_n}(t^*)|\sim u_0/n$. Integrating
\eqref{eq:continuity-fourier} from $0$ to $t^*$ with
$\hat\rho_{k_n}(0)\!\simeq\! 0$ (the step IC has flat density
initially) yields
\begin{equation}
|\hat\rho_{k_n}(t^*)| \sim k_n\,|\hat u_{k_n}^{(c)}(t^*)|\,t^*
\sim (2\pi/L)\,u_0\,t^*,
\label{eq:rhouratio}
\end{equation}
where $\hat u_{k_n}^{(c)}=L^{-1}\int u(s,t)e^{-ik_n s}ds$ is the
continuous Fourier coefficient (velocity units) and the
expression on the right uses
$|\hat u_{k_n}^{(c)}|\sim u_0/n$.
A $1/n$ cascade in velocity produces a density spectrum
$|\hat\rho_{k_n}|$ that is approximately $n$-independent, since
the carrier $k_n=2\pi n/L$ in Eq.~\eqref{eq:rhouratio} multiplies
$|\hat u_{k_n}|\sim 1/n$. The density-cascade signature is
therefore
\begin{equation}
\frac{|\hat\rho_{k_n}(t^*)|}{|\hat\rho_{k_1}(t^*)|}
\to \mathrm{const}\sim 1
\quad\text{(Burgers shock)}.
\end{equation}
In the weighted form
$n|\hat\rho_{k_n}|/|\hat\rho_{k_1}|$ used in the tables below
(retained for direct comparison with the velocity-cascade
exponent $1/n$), the Burgers prediction is therefore $\sim n$,
\emph{growing} with $n$. Note that this is a
\emph{linear-response} translation: at nonlinear amplitude,
$\widehat{(\delta\rho\,u)}_{k_n}$ provides additional mode
couplings whose phase is not constrained by continuity alone,
so the translation provides a definite signature only as long
as nonlinear corrections do not interfere destructively at
order unity. The diagnostic is sufficient (not necessary): a
true shock should leave it; absence of the growing-in-$n$
pattern is evidence against shock formation.

\subsection{Phase-mixing baseline}

In a non-interacting one-dimensional gas (and, by Jepsen's
mapping \cite{Jepsen1965}, in the pure Tonks gas) a thermal
sinusoidal density perturbation decays as a Gaussian
\cite{Lebowitz1968}
\begin{equation}
\hat\rho_k(t) = \hat\rho_k(0)\exp\!\left[-\tfrac{1}{2}(k c_s t)^2\right],
\quad c_s=\tfrac{v_{\rm th}}{1-\eta},
\label{eq:phasemix}
\end{equation}
without harmonic generation: the cascade ratio of
Eq.~\eqref{eq:burgersratio} stays at zero except for trivial
$O(A^2)$ corrections from the finite-amplitude initial condition.
The H\"ubner--Doyon theorem \cite{HubnerDoyon2024} implies that
no shock develops in this regime at any finite time, extending
the linear result to all orders in amplitude.

\subsection{Timescales: shock vs phase-mixing vs
integrability-breaking}\label{sec:timescales}

The dynamics admits three competing timescales relevant to shock
formation:
\begin{align}
t^*_{\rm shock}    &\sim 1/(k\,u_0)
                     \quad\text{(Burgers naive shock time)},
\label{eq:tshock}\\
\tau_G &= (1-\eta)\,L/(2\pi\,v_{\rm th})
         \quad\text{(phase-mixing)},
\label{eq:tauG}\\
\tau_{\rm break} &\sim 2(1-\eta)L/(N\bar v)
         \quad\text{(integrability-breaking)},
\label{eq:taubr}
\end{align}
with $\bar v=\sqrt{k_B T(m_1+m_2)/(m_1 m_2)}$.
Table~\ref{tab:timescales} reports these for the four
configurations of Sec.~\ref{sec:step}.

\begin{table}[t]
\centering
\caption{Competing timescales at the four step-IC
configurations ($L=12.57$, $N=400$, $\eta=0.1$, $k_B T=1$).
$t^*_{\rm shock}$ uses $u_0=U_0$, the maximum coherent velocity
of the step profile. The ordering
$\tau_{\rm break}\ll \tau_G\sim t^*_{\rm shock}$ holds in all
configurations: integrability-breaking is well established
before either phase-mixing or shock formation could complete,
so no-shock cannot be attributed to inadequate integrability
breaking or to phase-mixing pre-empting shock formation.}
\label{tab:timescales}
\begin{tabular}{cccc}
\hline\hline
config & $\tau_{\rm break}$ & $\tau_G$ & $t^*_{\rm shock}$\\
\hline
$m_5, U_{0.3}$   & 0.051 & 1.80 & 6.67 \\
$m_5, U_{0.5}$   & 0.051 & 1.80 & 4.00 \\
$m_{10},U_{0.3}$ & 0.054 & 1.80 & 6.67 \\
$m_{10},U_{0.5}$ & 0.054 & 1.80 & 4.00 \\
\hline\hline
\end{tabular}
\end{table}

Two consequences:
\emph{(i)} $\tau_{\rm break}\ll t^*_{\rm shock}$: the
mass-disordered dynamics has full opportunity to form a
shock before our observation window expires;
\emph{(ii)} $\tau_G\lesssim t^*_{\rm shock}$: the
phase-mixing window is comparable to or shorter than the
naive shock time, so we use this ordering to interpret
observations rather than excuse them. Phase-mixing must not
prevent us from detecting shock onset: the signature of a
forming shock is a transient $1/n$ cascade that should appear
\emph{before} $\tau_G$ has fully damped the signal.

\subsection{Burgers from mass-disordered hard rods?}

A naive two-fluid hydrodynamic argument suggests that binary
mass disorder should generate a viscous regularization of
Eq.~\eqref{eq:burgers}. The inter-species collision time
$\tau_{\rm coll}^{\rm inter}\sim
1/(\rho\sigma\langle|\Delta v|\rangle)$, combined with the
species' thermal velocity dispersion, yields a kinematic
viscosity
$\nu_{\rm kin}\sim v_{\rm th}\,\lambda_{\rm mfp}\approx 0.03$
at our parameters (with $\lambda_{\rm mfp}=(1-\eta)L/N$). The
expectation, then, is that step or sound-mode IC at sufficient
amplitude should develop viscous shocks described by
Eq.~\eqref{eq:burgers} with $\nu\sim\nu_{\rm kin}$.
Sec.~\ref{sec:phase} tests this expectation directly via the
Navier--Stokes residual and \emph{rejects} it.

\section{Tonks finite amplitude: HD2024 microscopically verified}
\label{sec:tonks}

We first establish the integrable baseline. Pure Tonks
($m_{\max}=1$) simulations across the three smooth IC families
and three amplitudes $A\in\{0.1,0.3,0.5\}$ show no harmonic
generation beyond the trivial $O(A^2)$ corrections.
Table~\ref{tab:tonks} summarises the maximum Burgers ratio
$n|\hat\rho_{k_n}|/|\hat\rho_{k_1}|$ over time.

\begin{table}[t]
\centering
\caption{Maximum weighted ratio
$n|\hat\rho_{k_n}|/|\hat\rho_{k_1}|$ in pure Tonks
(m$_{\max}$=1) finite-amplitude runs. A Burgers shock cascade
would give this weighted ratio $\sim n$, growing linearly; phase
mixing gives values decaying faster than constant. The ratios
observed are consistent with $O(A^2)$ nonlinear coupling only.
$^\dagger$Marked entries fall below the bootstrap signal-to-noise
threshold at the corresponding amplitude and are reported as
upper bounds.}
\label{tab:tonks}
\begin{tabular}{cccccc}
\hline\hline
IC & $A$ & $n=2$ & $n=3$ & $n=4$ & $n=5$\\
\hline
$\alpha$ & 0.10 & 0.05 & $<\!0.1^\dagger$ & $<\!0.1^\dagger$ & $<\!0.1^\dagger$\\
$\alpha$ & 0.30 & 0.02 & 0.02 & 0.04 & 0.04\\
$\alpha$ & 0.50 & 0.02 & 0.02 & 0.02 & 0.03\\
$\gamma$ & 0.30 & 0.11 & 0.03 & 0.03 & 0.05\\
$\gamma$ & 0.50 & 0.15 & 0.05 & 0.02 & 0.03\\
\hline\hline
\end{tabular}
\end{table}

Across all configurations, the ratio is bounded by
$\sim 0.15$ for $n=2$ and falls rapidly for $n\ge 3$. The
fundamental mode decays as a Gaussian with the predicted
timescale $\tau_G=(1-\eta)L/(2\pi v_{\rm th})$ to within
$\sim 2\%$. The pure Tonks gas does not develop Burgers
shocks at any of the amplitudes we probe, in agreement with
the H\"ubner--Doyon theorem.

\section{Mass-disordered finite amplitude: nonlinear sound
waves, no shock cascade}\label{sec:massdisorder}

We now break integrability with binary mass disorder
$m_i\in\{1,m_{\max}\}$. Companion paper~\cite{PaperII}
establishes that mass disorder restores sound waves in the
linear regime. We probe finite amplitude in the same
parameter space as Sec.~\ref{sec:tonks} at $m_{\max}=5$.
This $400\%$ mass disparity is far above the
$\sim 20\%$ at which linear-response sound waves first emerge
in companion paper~\cite{PaperII}; we are therefore well
within the broken-integrability regime in any reasonable
sense.

\begin{figure}[t]
\centering
\includegraphics[width=\columnwidth]{figures/p3_secIV_logmodes.pdf}
\caption{Density Fourier-mode amplitudes $|\hat\rho_{k_n}|$ for
$n=1\ldots 5$ in the mass-disordered hard-rod gas ($m_{\max}=5$,
$N=400$, $\eta=0.10$) with density-only IC at $A=0.5$.
Top: time evolution on a logarithmic vertical scale.
The fundamental oscillates at the sound-wave frequency and decays;
higher harmonics remain suppressed, with no Burgers cascade.
Bottom: peak weighted ratio $n|\hat\rho_{k_n}|/|\hat\rho_{k_1}|$,
which decays with $n$ rather than growing as a Burgers shock
would require.}
\label{fig:massdisorder}
\end{figure}

Figure~\ref{fig:massdisorder} (top) shows the time evolution of
$|\hat\rho_{k_n}|$ for $n=1\ldots 5$ in the density-only IC at
$A=0.5$. The fundamental mode oscillates with the
sound-wave frequency of \cite{PaperII} and decays slowly. The
second harmonic appears with amplitude $\sim 0.19$, an order of
magnitude smaller than $|\hat\rho_{k_1}|$, and the higher
harmonics remain at the percent level. The peak Burgers ratio
$n|\hat\rho_{k_n}|/|\hat\rho_{k_1}|$ versus $A$
(Table~\ref{tab:msd}) increases with amplitude but stays well
below unity, and---crucially---\emph{decreases} with $n$ rather
than remaining constant. The implied cascade follows $1/n^2$
(smooth wave) rather than $1/n$ (Burgers shock).

\begin{table}[t]
\centering
\caption{Weighted ratio $n|\hat\rho_{k_n}|/|\hat\rho_{k_1}|$
in mass-disordered ($m_{\max}=5$) runs. Burgers shock prediction
$\sim n$ (growing). The ratio increases with $A$ but is bounded
by $\sim 0.5$ even at $A=0.8$, and \emph{decreases} sharply with
$n$, opposite to the Burgers prediction.}
\label{tab:msd}
\begin{tabular}{ccccc}
\hline\hline
IC, $A$ & $n=2$ & $n=3$ & $n=4$ & $n=5$\\
\hline
$\alpha$, 0.1 & 0.05 & 0.02 & 0.01 & 0.01\\
$\alpha$, 0.3 & 0.24 & 0.06 & 0.02 & 0.02\\
$\alpha$, 0.5 & 0.39 & 0.13 & 0.07 & 0.04\\
$\alpha$, 0.8 & 0.52 & 0.30 & 0.16 & 0.11\\
\hline\hline
\end{tabular}
\end{table}

Runs at sound-mode IC ($\gamma$) and at $m_{\max}=10$ are
qualitatively similar but are not tabulated in detail here; the
$n=2$ ratio peaks at $0.41$ ($\gamma$, $A=0.5$, $m_{\max}=5$) and
$0.58$ ($\alpha$, $A=0.8$, $m_{\max}=10$), with $n\geq 3$
suppressed in both cases. The system is in a \emph{weakly
nonlinear sound-wave regime} with $O(A^2)$ mode coupling, not
in a shock-forming regime.

\section{Step velocity IC: forced shock attempt also fails}
\label{sec:step}

A natural concern is that the smooth sinusoidal IC simply does
not develop sharp enough gradients to trigger shocks within our
simulation window. We address this with a step velocity IC,
$u_{\rm coh}(s)=U_0\,\mathrm{sgn}[\cos(ks)]$, the canonical
Riemann initial data for Burgers. The initial Fourier content
is $|\hat u_{k_n}(0)|=(2U_0/\pi)/n$ for odd $n$ in our complex
convention (Sec.~\ref{sec:setup}), already following the Burgers
$1/n$ velocity-cascade exponent. If the
dynamics is governed by an inviscid Burgers equation, this
profile would be a near-stationary state.

\begin{figure}[t]
\centering
\includegraphics[width=\columnwidth]{figures/p3_stepUB_spectrum.pdf}
\caption{Density-mode spectrum at the peak time $t^*$ (maximum of
$|\hat\rho_{k_2}|$) for step velocity IC with
$m_{\max}\in\{5,10\}$, $U_0=0.5$, $N=400$, $\eta=0.10$.
The weighted ratio $n|\hat\rho_{k_n}|/|\hat\rho_{k_1}|$ shows a
transient $n=2$ enhancement from second-order nonlinear coupling
of odd modes, but no sustained cascade: the ratio drops sharply
for $n\ge 3$, inconsistent with the Burgers prediction $\sim n$.}
\label{fig:step}
\end{figure}

Figure~\ref{fig:step} shows the density-mode spectrum at the
peak time $t^*$ (defined as the maximum of
$|\hat\rho_{k_2}(t)|$) for $m_{\max}\in\{5,10\}$,
$U_0=0.5$. Two features stand out, both deviating from the
Burgers prediction $n|\hat\rho_{k_n}|/|\hat\rho_{k_1}|\sim n$:

\noindent
(i) the ratio at $n=2$ takes values $2.06$ ($m_{\max}=5$) and
$4.08$ ($m_{\max}=10$), within a factor of $2$ of the Burgers
prediction $n=2$. However, the step initial condition has no
initial content at $n=2$ (only odd modes populated), and we
show below that this enhancement is generated by second-order
nonlinear coupling between linearly-populated odd modes, with
no extension to a higher-order cascade;

\noindent
(ii) for $n\ge 3$, the ratio drops to $0.2$--$0.5$, an order
of magnitude below the Burgers prediction $n=3,4,5,\ldots$
and decaying rather than growing with $n$.

Table~\ref{tab:step} reports the cascade ratios for the four
step-IC configurations.

\begin{table}[t]
\centering
\caption{Weighted ratio $n|\hat\rho_{k_n}|/|\hat\rho_{k_1}|$ in
step-IC runs at the peak time of $|\hat\rho_{k_2}|$.
Burgers shock prediction $\sim n$ (growing linearly). Observed
ratios decay with $n$ from $n=3$ onwards, inconsistent with a
sustained cascade.}
\label{tab:step}
\begin{tabular}{cccccccc}
\hline\hline
config & $n=1$ & $n=2$ & $n=3$ & $n=4$ & $n=5$ & $n=6$ & $n=7$\\
\hline
$m_5, U_{0.3}$  & 1.0 & 0.74 & 0.05 & 0.05 & 0.09 & 0.08 & 0.06\\
$m_5, U_{0.5}$  & 1.0 & 2.06 & 0.35 & 0.26 & 0.44 & 0.34 & 0.21\\
$m_{10},U_{0.3}$& 1.0 & 0.63 & 0.05 & 0.10 & 0.16 & 0.16 & 0.16\\
$m_{10},U_{0.5}$& 1.0 & 4.08 & 0.25 & 0.53 & 1.23 & 0.67 & 0.53\\
\hline\hline
\end{tabular}
\end{table}

\subsection{Supersonic regime: same result, larger amplitudes}
\label{sec:supersonic}

A possible objection is that $U_0=0.5$ is subsonic
($U_0/c_s\approx 0.45$, $c_s = v_{\rm th}/(1-\eta)\approx 1.11$),
so the system has not been tested in the regime where Burgers
nonlinearity is naively dominant. We extend the step-IC scan to
$U_0\in\{1.0, 1.5, 2.0\}$ at $m_{\max}=5$, the upper end
corresponding to $U_0/c_s\approx 1.80$.
Table~\ref{tab:super} reports the dimensionless density-mode
ratios $\max_t|\hat\rho_{k_n}|/\max_t|\hat\rho_{k_2}|$ (use of
the max-over-time normalisation avoids artefacts from the
oscillation of $|\hat\rho_{k_1}|$ near nodes).

\begin{table}[t]
\centering
\caption{Density-mode ratios from supersonic step-IC scan,
$m_{\max}=5$, $4000$ seeds. The shape of the spectrum is
essentially independent of $U_0/c_s$ across the subsonic-to-%
supersonic transition; modes scale uniformly with amplitude
while no Burgers cascade emerges.}
\label{tab:super}
\begin{tabular}{cccccccc}
\hline\hline
$U_0$ & $U_0/c_s$ & $|\hat\rho_1|/|\hat\rho_2|$
       & $|\hat\rho_3|/|\hat\rho_2|$
       & $|\hat\rho_4|/|\hat\rho_2|$
       & $|\hat\rho_5|/|\hat\rho_2|$
       & $|\hat\rho_6|/|\hat\rho_2|$
       & $|\hat\rho_7|/|\hat\rho_2|$\\
\hline
$0.5$ & $0.45$ & $3.83$ & $1.28$ & $0.39$ & $0.77$ & $0.25$ & $0.54$\\
$1.0$ & $0.90$ & $3.24$ & $1.08$ & $0.39$ & $0.64$ & $0.24$ & $0.46$\\
$1.5$ & $1.35$ & $3.16$ & $1.05$ & $0.37$ & $0.63$ & $0.23$ & $0.44$\\
$2.0$ & $1.80$ & $3.18$ & $1.05$ & $0.41$ & $0.61$ & $0.25$ & $0.45$\\
\hline\hline
\end{tabular}
\end{table}

Three observations support the no-shock conclusion at
supersonic amplitudes:
\emph{(i)} the spectral ratios are essentially $U_0$-independent
across the sub-to-supersonic transition;
\emph{(ii)} the odd-$n$ modes ($n=3,5,7$) decay with $n$
(values $1.05, 0.63, 0.45$ at $U_0/c_s=1.80$) rather than
remaining near unity as a preserved Burgers cascade would
require;
\emph{(iii)} the even-$n$ modes ($n=4,6$), which the step IC
does not populate initially, remain suppressed by a factor of
$2$--$3$ relative to the adjacent odd modes, signalling weak
nonlinear coupling rather than a shock-mediated cascade.

We conclude that the absence of Burgers shock formation persists
into the supersonic regime accessible to our simulations.

\subsection{Heavy-mass limit: same result}
\label{sec:heavy}

A complementary regime is $m_{\max}\to\infty$, where heavy rods
become effectively immobile on the simulation timescale
($v_{\rm th, H}=1/\sqrt{m_{\max}}\ll U_0$) and the system
approaches a Tonks gas of light rods scattering off a fixed
random array of impurities. Naively one expects this to
provide an effective viscous regularisation and possibly
restore Burgers shocks. We test $m_{\max}\in\{20, 100\}$
(step IC, $U_0=0.5$, $4000$ seeds each).

\begin{table}[t]
\centering
\caption{Density-mode ratios $\max_t|\hat\rho_{k_n}|
/\max_t|\hat\rho_{k_2}|$ in the heavy-mass limit. Across two
orders of magnitude in $m_{\max}$ the spectral shape is
essentially unchanged, ruling out the alternative that the
no-shock result is an artifact of weakly-disordered dynamics.}
\label{tab:heavy}
\begin{tabular}{ccccccc}
\hline\hline
$m_{\max}$ & $|\hat\rho_1|/|\hat\rho_2|$
           & $|\hat\rho_3|/|\hat\rho_2|$
           & $|\hat\rho_4|/|\hat\rho_2|$
           & $|\hat\rho_5|/|\hat\rho_2|$
           & $|\hat\rho_6|/|\hat\rho_2|$
           & $|\hat\rho_7|/|\hat\rho_2|$\\
\hline
$5$   & $3.84$ & $1.28$ & $0.37$ & $0.76$ & $0.24$ & $0.55$\\
$20$  & $3.21$ & $1.06$ & $0.39$ & $0.64$ & $0.23$ & $0.44$\\
$100$ & $3.61$ & $1.19$ & $0.46$ & $0.70$ & $0.26$ & $0.49$\\
\hline\hline
\end{tabular}
\end{table}

The density-mode spectrum (Table~\ref{tab:heavy}) is statistically
indistinguishable across $m_{\max}=5,20,100$ in normalised form.
The absolute amplitudes scale up modestly with $m_{\max}$
(heavy rods reflect light rods more elastically, transferring
more amplitude to the density field), but the spectral shape
--- the diagnostic for shock formation --- is unaffected. The
no-shock result is therefore robust across the full
mass-disorder parameter range, from weak ($m_{\max}=5$) to
effectively static heavy-rod arrays ($m_{\max}=100$).

\subsection{Origin of the $n=2$ enhancement}

The step IC carries only odd-$n$ velocity Fourier content:
direct evaluation of
$\hat u^{(c)}_n=L^{-1}\!\int U_0\,\mathrm{sgn}[\cos(ks)]\,
e^{-ik_n s}\,ds$ gives magnitudes $2U_0/(\pi n)$ for odd $n$
and zero for even $n$. Linearised continuity in Fourier space,
$\partial_t \hat\rho_n = -ik_n \hat u^{(c)}_n$, then populates
only odd density modes at first order in $U_0$. The $n=2$
enhancement is therefore \emph{not} a linear-kinematic
feature; it arises from second-order nonlinear coupling between
the linearly-generated odd modes via the
$\widehat{(\delta\rho\,u)}$ term in continuity. A naive
second-order estimate is
$|\hat\rho_{k_2}(t^*)|^{(2)} \sim k_1 k_2 t^{*2}
|\hat u_{k_1}^{(c)}|^2 \approx 0.5$
for $U_0=0.5$, $t^*\approx 3.4$, $L=12.57$. The observed
$|\hat\rho_{k_2}(t^*)| \approx 0.068$ is a factor of seven
below this second-order estimate, indicating that the coupling
is itself suppressed by subsequent collisional mixing.

The diagnostic that discriminates is the cascade at $n\ge 3$.
A true Burgers shock developing from the step IC would preserve
the initial $1/n$ velocity-cascade exponent (Sec.~\ref{sec:rhoVSu}),
giving weighted ratios $n|\hat\rho_{k_n}|/|\hat\rho_{k_1}|
\sim 3, 4, 5,\ldots$ (growing linearly with $n$). The observed
ratios at $n\ge 3$ are $0.2$--$0.5$ (Table~\ref{tab:step}), an
order of magnitude or more below the Burgers prediction and
\emph{decaying} with $n$ rather than growing. We read this as
evidence that no genuine Burgers cascade develops.

\section{Arc-length universality at finite amplitude}
\label{sec:curved}

The arc-length reduction theorem of companion paper~\cite{PaperI}
states that, for a smooth closed curve $\gamma$ in
$\mathbb{R}^d$ with arc-length element
$ds=\sqrt{g(\varphi)}d\varphi$, the hard-rod dynamics is
\emph{exactly} that of the flat Tonks gas of perimeter $L$. The
theorem is established in the linear-response regime by
explicit computation of equilibrium correlators
\cite{PaperI,PaperII}. Here we verify that the conclusion
extends to the strongly nonlinear, step-IC regime studied
above.

We repeat the step-IC simulation ($m_{\max}=5$, $U_0=0.5$) on
the wavy embedding $r(\varphi)=R[1+0.3\cos(3\varphi)]$ with
$L=14.85$, $44\%$ longer than the circular reference
$L=12.57$. Coincidence with the circle is expected on
theoretical grounds; the data
(Table~\ref{tab:curved}) confirm it within statistical error.

\begin{table}[t]
\centering
\caption{Comparison of cascade ratios on circle and on the
wavy ring $r=R[1+0.3\cos(3\varphi)]$ at matched $m_{\max}=5$,
$U_0=0.5$. The metric profile changes substantially while the
result is unaffected.}
\label{tab:curved}
\begin{tabular}{ccccccc}
\hline\hline
geometry & $n=2$ & $n=3$ & $n=4$ & $n=5$ & $n=6$ & $n=7$\\
\hline
circle ($L=12.57$) & 2.06 & 0.35 & 0.26 & 0.44 & 0.34 & 0.21\\
wavy ($L=14.85$)   & 2.04 & 0.33 & 0.21 & 0.40 & 0.42 & 0.26\\
\hline\hline
\end{tabular}
\end{table}

This constitutes a nonlinear-regime verification of the
arc-length reduction theorem: the density-mode spectrum
depends on the geometry only through $L$, even at finite
amplitude with a metric strongly varying in $\varphi$.
The result is expected; reporting it here closes the loop with
Refs.~\cite{PaperI,PaperII} and confirms that the no-shock
conclusion of Secs.~\ref{sec:tonks}--\ref{sec:step} is a property
of the flat-Tonks (binary-mixture) dynamics, not an artefact
of the chosen geometry.

\section{Phase-space moment analysis}
\label{sec:phase}

To diagnose the failure of Burgers more directly, we compute
spatially-binned phase-space moments
\begin{align}
\rho(s,t) &= \frac{1}{\Delta s}\,\langle\#\{i:s_i(t)\in
\mathrm{bin}_s\}\rangle, \\
u(s,t) &= \langle v_i\rangle_{\mathrm{bin}_s}, \\
T(s,t) &= \langle (v_i - u)^2\rangle_{\mathrm{bin}_s},
\end{align}
on $80$ uniform bins of width $\Delta s = L/80$, averaged over
$3000$ realisations. We then evaluate the residuals of the
continuity and Burgers equations:
\begin{align}
R_{\rm cont}(s,t)&=\partial_t \rho + \partial_s (\rho u),\\
R_{\rm Burg}(s,t)&=\partial_t (\rho u) + \partial_s (\rho u^2 + P),
\end{align}
with the \emph{local} Tonks pressure
$P(s,t) = \rho(s,t)\,T(s,t) /[1-\sigma\rho(s,t)]$. Continuity is an identity of the microscopic dynamics, holding
in the joint limit $M\to\infty$, $\Delta s\to 0$. At finite
$M$, the bin-averaged $\rho$ and $\rho u$ acquire shot-noise of
amplitude $\sim 1/\sqrt{M \rho_0 \Delta s}$, which propagates to
$R_{\rm cont}$ at a level $\sim 0.06$ in our $M=6000$ runs at
$n_{\rm bins}=80$ --- consistent with the value reported in
Table~\ref{tab:refine}. The Burgers residual is a test of
second-moment closure with this specific pressure choice. We have verified that the simpler global-$\eta$ pressure
$P=\rho T/(1-\eta)$ gives statistically indistinguishable
residuals for the step IC (Table~\ref{tab:refine}), because
$\sigma\rho(s,t)$ stays close to $\eta$ throughout the
simulation window: a velocity step does not, by itself, generate
order-unity density variations on the time scales accessed.

We compute the peak ratio of each residual to the local rate of
change of the corresponding signal,
$\max|R_{\rm cont}|\,\Delta t / \max|\Delta\rho|$ and
$\max|R_{\rm Burg}|\,\Delta t / \max|\Delta(\rho u)|$.
Table~\ref{tab:moments} reports these for the three IC families
at $A=0.5$.

\subsubsection*{Bin convergence and bootstrap errors}

We performed a dedicated refinement run (step IC, $m_{\max}=5$,
$U_0=0.5$, $6000$ seeds organised in $30$ chunks of $200$ for
chunk-bootstrap) at three spatial resolutions
$n_{\rm bins}\in\{40,80,160\}$.
Table~\ref{tab:refine} reports point estimates and $1\sigma$
bootstrap errors for the continuity residual and for the Burgers
residual under both the global-$\eta$ and local Tonks pressures.

\begin{table}[t]
\centering
\caption{Continuity and Burgers residuals at three spatial
resolutions, step IC, $m_{\max}=5$, $U_0=0.5$.
$1\sigma$ chunk-bootstrap errors quoted ($1000$ resamples on
$30$ chunks of $200$ seeds). The Burgers residual grows
between $n_{\rm bins}=40$ and $80$ from $0.40$ to $0.76$ while
continuity stays at the per-bin shot-noise floor; at
$n_{\rm bins}=160$ shot noise dominates both residuals
(particle count per bin $\approx 2.5$, below the regime where
finite differences are meaningful). Global-$\eta$ and local
$\sigma\rho$ pressures yield statistically indistinguishable
results at $n_{\rm bins}\le 80$.}
\label{tab:refine}
\begin{tabular}{cccc}
\hline\hline
$n_{\rm bins}$ & continuity & Burgers $P(\eta)$
                                  & Burgers $P(\sigma\rho)$\\
\hline
$40$  & $0.070\pm 0.003$ & $0.403\pm 0.008$ & $0.403\pm 0.008$\\
$80$  & $0.056\pm 0.003$ & $0.759\pm 0.014$ & $0.759\pm 0.014$\\
$160$ & $0.881\pm 0.011$ & $0.838\pm 0.024$ & $0.949\pm 0.026$\\
\hline\hline
\end{tabular}
\end{table}

Two facts isolate a genuine closure-failure signal from
discretisation:
\emph{(i)} the local and global-$\eta$ pressures give
indistinguishable Burgers residuals at $n_{\rm bins}=40, 80$
(the regime where shot noise has not yet dominated). This rules
out the alternative that the residual is an artifact of the
mean-field pressure approximation.
\emph{(ii)} continuity, an exact identity, satisfies
$R_{\rm cont}\sim 0.06$--$0.07$ at $n_{\rm bins}\le 80$ while
the Burgers residual at the same resolutions is one order of
magnitude larger. Both quantities use identical finite
differences, so this is not a discretisation effect.

\subsubsection*{Local thermodynamic equilibrium}

A second concern is whether the local velocity distribution is
sufficiently equilibrated to justify using an equilibrium Tonks
pressure at all: for a step velocity IC, the local distribution
could in principle be a bimodal two-stream, in which case
$\langle (v-u)^2\rangle$ is not a true temperature and
$P=\rho T/(1-\sigma\rho)$ is not the correct momentum flux. We
test this by computing the fourth central moment per bin and
the kurtosis $\kappa = \langle(v-u)^4\rangle/T^2$. A bimodal
two-stream would yield $\kappa\to 1$; Maxwell--Boltzmann
$\kappa = 3$; our data give a median $\kappa\approx 4.33$ across
all space-time bins with $q_{95}\approx 4.69$, and the global
minimum is $\kappa_{\min}\approx 3.76$ --- nowhere does the
distribution become bimodal. The deviation from
$\kappa = 3$ is explained quantitatively by the binary-mass
mixture: a $50$/$50$ mixture of Gaussians with
$\sigma_v\in\{1,1/\sqrt{m_{\max}}\}$ at $m_{\max}=5$ has
$\kappa_{\rm mix} = 4.33$, in numerical agreement with the
observed value. Since $T=\langle(v-u)^2\rangle$ is by definition
the correct ideal-gas momentum flux per particle regardless of
the underlying distribution, the pressure expression we use
remains physically appropriate; the kurtosis result indicates
that the gas is in local equilibrium with the binary-mass
distribution, not in a bimodal non-equilibrium state.

\subsubsection*{Irreversibility without shocks}

A natural worry is that the system might simply lack
dissipation, in which case ``no shock'' would be unremarkable.
We diagnose this via the partition of total kinetic energy
into coherent and thermal components,
$E_{\rm coh}(t) = \tfrac12\sum_b (\rho_b^L + m_{\max}\rho_b^H)\,
u_b^2$
and
$E_{\rm th}(t) = E_{\rm tot}-E_{\rm coh}$, with $u_b$ the
mass-weighted mean velocity in bin $b$ and species-resolved
densities $\rho_b^{L,H}$. Total energy is conserved to numerical
precision ($\Delta E/E < 10^{-3}$). The coherent component
decays rapidly: from a step-IC initial fraction
$E_{\rm coh}(0)/E_{\rm tot}=0.43$, half-life is
$t_{1/2}\approx 1.2$ and $E_{\rm coh}/E_{\rm tot}<0.10$ by
$t\approx 2.2$, well before the naive Burgers shock time
$t^*_{\rm shock}\sim 4$ (Table~\ref{tab:ent}).

\begin{table}[t]
\centering
\caption{Coherent vs thermal kinetic-energy partition, step IC,
$m_{\max}=5$, $U_0=0.5$. Total energy is exactly conserved by
the elastic-collision dynamics; the coherent fraction decays
on a time scale of $\sim 1$, much faster than the naive
Burgers shock time $t^*_{\rm shock}\approx 4$.}
\label{tab:ent}
\begin{tabular}{cccc}
\hline\hline
$t$ & $E_{\rm coh}$ & $E_{\rm th}$ & $E_{\rm coh}/E_{\rm tot}$\\
\hline
$0.0$  & $149.6$ & $198.9$ & $0.429$\\
$1.0$  & $\phantom{0}90.6$ & $258.0$ & $0.260$\\
$3.0$  & $\phantom{00}5.5$ & $343.0$ & $0.016$\\
$10.0$ & $\phantom{0}27.8$ & $320.8$ & $0.080$\\
$30.0$ & $\phantom{00}0.4$ & $348.2$ & $0.001$\\
\hline\hline
\end{tabular}
\end{table}

Kinetic equipartition between species is preserved to better
than $1\%$ throughout the simulation: $\tfrac12 m_L\langle
v^2\rangle_L = \tfrac12 m_H\langle v^2\rangle_H = \tfrac12 k_B T$,
giving $T_H/T_L = 1/m_{\max}$ exactly. This is a non-trivial
observation: in a system that has only just broken
integrability, equipartition is a thermalisation result, not
a kinematic identity, and its preservation in the presence of
a strong coherent perturbation shows that inter-species
thermalisation is faster than the coherent-mode dynamics.
The interpretation is therefore unambiguous: the mass-disordered
gas \emph{does} dissipate coherent kinetic energy on
coarse-grained scales --- the underlying elastic dynamics is
strictly time-reversible, but the projection onto the
hydrodynamic variables $\{\rho, u, T\}$ produces monotonic decay
of $E_{\rm coh}$ via inter-species collisions, on a time scale
of order $\tau_{\rm break}$. Yet no Burgers shock forms. The
``no-shock'' result of Sec.~\ref{sec:step} cannot be attributed
to absent dissipation at the coarse-grained level; it reflects
the fact that this dissipation channel is realised as
phase-mixing rather than as shock-mediated gradient sharpening.

\subsubsection*{Viscous closure does not rescue Euler}

A natural alternative reading of a nonzero Burgers residual is
that we have omitted the viscous correction
$\partial_s(\nu\partial_s u)$ from a Navier--Stokes-type
closure. To test this, we add the viscous term to the residual
and sweep $\nu\in[0,100]$, recording the minimum normalised
residual $\min_\nu R_{\rm NS}(\nu)$ and the optimal $\nu^*$
(Table~\ref{tab:nu}). Three observations disfavour the viscous
interpretation:
\emph{(i)} the minimum residual remains an order of magnitude
above the continuity floor;
\emph{(ii)} the optimal $\nu^*\sim 3$--$8$ is two orders of
magnitude larger than the kinetic-theory estimate
$\nu_{\rm kin}\sim v_{\rm th}\lambda_{\rm mfp}\approx 0.03$
based on the inter-species mean free path;
\emph{(iii)} $\nu^*$ depends on the spatial resolution, which a
genuine transport coefficient should not. Navier--Stokes
closure with a uniform shear viscosity does not, therefore,
account for the residual.

\begin{table}[t]
\centering
\caption{Viscous-Navier--Stokes closure test: minimum residual
$\min_\nu R_{\rm NS}$ and optimal $\nu^*$ from a sweep
$\nu\in[0,100]$, step IC. Even with the optimal $\nu$, the
residual remains far above the continuity floor; $\nu^*$ is
two orders of magnitude above the kinetic estimate and
resolution-dependent, ruling out a viscous interpretation.}
\label{tab:nu}
\begin{tabular}{cccc}
\hline\hline
$n_{\rm bins}$ & $R_{\rm Euler}$ ($\nu=0$)
                  & $\min_\nu R_{\rm NS}$ & $\nu^*$\\
\hline
$40$  & $0.403$ & $0.280$ & $8.0$\\
$80$  & $0.759$ & $0.515$ & $4.0$\\
$160$ & $0.949$ & $0.865$ & $3.5$\\
\hline\hline
\end{tabular}
\end{table}

\begin{table}[t]
\centering
\caption{Residuals across initial-condition families at the
working resolution $n_{\rm bins}=80$, $m_{\max}=5$ on circle.
For smooth IC ($\alpha$, $\gamma$) the Burgers residual is
comparable to continuity (both within the discretisation
floor): no closure failure detectable. For the step IC ($\delta$)
the Burgers residual is $\sim 10\times$ larger than continuity
at identical discretisation, isolating a genuine
sharp-gradient closure failure.}
\label{tab:moments}
\begin{tabular}{cccc}
\hline\hline
IC & continuity & Burgers ($P_{\rm Tonks}$ local) \\
\hline
$\alpha$ ($A=0.5$)        & $0.28$ & $0.32$ \\
$\delta$ (step, $U_0=0.5$)& $0.06$ & $\mathbf{0.76}$ \\
$\gamma$ ($A=0.5$)        & $0.24$ & $0.18$ \\
\hline\hline
\end{tabular}
\end{table}

Combining Tables~\ref{tab:moments} and \ref{tab:refine}, the
reading is: for smooth initial conditions the Burgers
residual is within the discretisation floor (no Euler closure
failure detectable); for the sharp-gradient step IC the
Burgers residual exceeds the floor by an order of magnitude at
$n_{\rm bins}=80$ and grows with resolution up to where shot
noise sets in. The second-moment closure of the
mass-disordered hard-rod gas \emph{fails} for sharp velocity
gradients. This is consistent with the natural reading from
generalised hydrodynamics: the relevant continuum variables
are not $\{\rho, u, T\}$ but the full rapidity-resolved
quasiparticle density.

\section{Discussion}\label{sec:discussion}

\subsection{Relation to H\"ubner--Doyon}

The H\"ubner--Doyon theorem~\cite{HubnerDoyon2024} establishes
analytically that the integrable Lieb--Liniger gas does not
develop Burgers shocks at any finite time, exploiting a new
quadrature of the GHD equation. The Lieb--Liniger and the
classical hard-rod gas are distinct models but share a common
theme: both are integrable, both have hydrodynamic limits
governed by GHD, and both exhibit the same absence-of-shocks
phenomenology. Our microscopic verification in the
hard-rod case (Sec.~\ref{sec:tonks}) confirms the classical
analogue of the H\"ubner--Doyon result and serves as a baseline
against which the mass-disorder results
(Sec.~\ref{sec:massdisorder}--\ref{sec:heavy}) are
compared.

\subsection{Why mass disorder may not be enough}

The data of Sec.~\ref{sec:phase} show that the no-shock result
is not for want of dissipation: total energy is exactly
conserved while coherent kinetic energy decays on a time scale
$\sim 1$ into thermal motion (Table~\ref{tab:ent}). Yet this
dissipation does not produce shock-cascade Fourier
spectra. The conserved quantities surviving the mass-disorder
breaking are limited: the global $N_1,N_2,P_{\rm tot},
E_{\rm tot}$, plus the structural fact that the gas remains a
sum of two conserved-mass species coupled only by point
elastic collisions. We do not claim a residual quasi-integrable
structure beyond these; the question of which (if any) further
generalised conservation laws constrain the dynamics on
intermediate time scales is open. The fact that the
$m_{\max}\to\infty$ limit gives a spectrum
statistically indistinguishable from $m_{\max}=5$
(Sec.~\ref{sec:heavy}) is itself a non-trivial constraint on
any candidate effective description: heavy-mass mixtures and
weakly-disordered mixtures share the same large-scale
hydrodynamics.

\subsection{Mass-disorder universality of the density spectrum}

The most surprising single observation of this study is
captured in Table~\ref{tab:heavy}: the density-mode spectrum
shape $|\hat\rho_{k_n}|/|\hat\rho_{k_2}|$ is statistically
indistinguishable across $m_{\max}\in\{5,20,100\}$ in the
step-IC dynamics --- a 20-fold range in the
integrability-breaking strength. This is not a trivial scaling:
the absolute amplitudes do grow with $m_{\max}$ (by
factor of two between $m_{\max}=5$ and $m_{\max}=100$),
so the system is genuinely sensitive to the impurity mass. What
collapses is the \emph{spectral shape}: the relative weight of
each Fourier mode in the density-mode response is set by
something other than $m_{\max}$. We do not have a closed
explanation for this universality and we present it as an open
challenge: any candidate effective theory for the
mass-disordered hard-rod gas must reproduce a density-mode
spectrum whose dimensionless shape is independent of $m_{\max}$
over at least an order of magnitude.

\subsection{Outlook}

In the GHD language, the natural candidate continuum framework
is the GHD equation extended to two species coupled by an
inter-species transfer kernel of strength
$\sim 1/\tau_{\rm break}$. We do not develop this equation
here, and we do not claim that GHD plus a transfer kernel is
the correct effective theory: our positive results are
the negative observations on shock formation and on
second-moment closure failure. The construction of a perturbed
two-species GHD equation and its quantitative test against the
data is a natural follow-up but lies outside the present
scope.

\subsection{What ingredient is missing?}

Lighthill--Whitham--Richards models of vehicular traffic
\cite{Lighthill1955} produce shocks at the continuum level
from microscopic heterogeneity alone. Our result indicates
that in a clean classical analogue (elastic mass-disordered
hard rods) this expectation is not met; reaction delay,
anticipation horizon, or genuinely dissipative collisions would
each be candidate additional ingredients, but testing them is
outside the present scope.

\subsection{Relevance for cold-atom GHD experiments}

For the cold-atom realisations of one-dimensional gases in
which generalised hydrodynamics has been quantitatively tested
\cite{Schemmer2019,Malvania2021,Cataldini2022}, our results have
a concrete operational implication: a controlled mass-disparity
impurity --- a routine experimental knob in mixture experiments
--- does not by itself displace the dynamics from the GHD
basin, even at strongly nonlinear perturbation amplitudes.
Finite-amplitude shock-cascade phenomenology would therefore
require a different microscopic ingredient, presumably one that
breaks the elastic two-body collision structure entirely.

\subsection{Settling the open question}

The cleanest theoretical test of the GHD interpretation would
be a direct two-species GHD equation solved numerically against
our microscopic data: a positive comparison would close the
chapter, a negative comparison would point to a genuinely
non-GHD effective theory. On the experimental side, a
fluctuation-theorem-style measurement of the work
distribution during the step-IC relaxation would resolve
whether the coarse-grained irreversibility we report
(Sec.~\ref{sec:phase}) carries a thermodynamic-entropy
signature beyond the energy-partition diagnostic. Both lines
of work are natural follow-ups.

\section{Conclusion and open work}\label{sec:conclusion}

We have shown that, on the time scales and amplitudes accessible
to our event-driven molecular dynamics (including the
supersonic regime $U_0/c_s \approx 1.8$ where Burgers
nonlinearity is naively dominant), binary mass disorder is
\emph{not} sufficient to restore Burgers shocks in the
hard-rod gas, despite breaking the velocity-multiset
conservation of the integrable Tonks limit on a time scale
$\tau_{\rm break}\sim 0.05$, two orders of magnitude shorter
than the naive shock time. The mode-$n$ density spectrum
decays faster than $1/n$ across all four initial-condition
families, four amplitudes and two mass disparities tested; the
arc-length universality of the companion paper carries over to
the strongly nonlinear regime; and a phase-space moment
analysis with bin-convergence study and chunk-bootstrap errors
demonstrates that the second-moment Euler closure
\emph{fails} for sharp velocity gradients under both
global-$\eta$ and local Tonks pressures, with the Burgers
residual one order of magnitude above the shot-noise floor of
the continuity residual.

No second-moment closure --- Euler with global-$\eta$ or local
Tonks pressure, or Navier--Stokes with uniform viscosity ---
reproduces the dynamics for sharp gradients. The candidate
effective theory remains the generalised hydrodynamic equation,
which we do not develop here. Genuine shock formation in
one-dimensional classical gases plausibly requires a
microscopic ingredient --- dissipation, finite reaction time, or
inelastic collisions --- beyond elastic mass disorder. We list as open theoretical work:
\emph{(i)} a direct test of whether a perturbed GHD equation
with an inter-species transfer kernel can reproduce the
observed density-mode evolution;
\emph{(ii)} extension to dissipative collisions or
finite-reaction-time microdynamics as candidate ingredients for
genuine shock formation;
\emph{(iii)} bootstrap error bars on the Fourier-mode tables
(Tables~\ref{tab:tonks}--\ref{tab:curved}), straightforward
extensions of the chunk-bootstrap performed for
Table~\ref{tab:refine}.

\begin{acknowledgments}
This work was supported by Tecnol\'ogico de Monterrey.
Computations were performed on the xolotl cluster.
\end{acknowledgments}

\appendix

\section{Unbiased Fourier-mode estimator}\label{app:estimator}

Computing the magnitude $|\hat\rho_{k_n}|=\sqrt{a_n^2+b_n^2}$
per realisation and averaging over realisations introduces a
positive bias of order $1/\sqrt{N}\approx 0.05$ at $N=400$. The
unbiased estimator separates the cosine and sine components,
$a_n^{(s)}=N^{-1}\sum_i \cos(k_n s_i^{(s)})$ and analogously
$b_n^{(s)}$ per realisation $s$, averages each over
realisations, $\bar a_n = \langle a_n^{(s)}\rangle$,
$\bar b_n = \langle b_n^{(s)}\rangle$, and recovers the
magnitude $|\hat\rho_{k_n}| = \sqrt{\bar a_n^2 + \bar b_n^2}$.
The bootstrap standard error scales as
$1/\sqrt{N M}$ for $M$ realisations, giving $\sim 8\times 10^{-4}$
in our highest-statistics runs (\textsc{stepUB}, $M=4000$).

\section{Phase-space moment computation}\label{app:moments}

For each realisation $s$, particles are binned in $80$ uniform
$s$-bins, and for each bin the per-realisation count $c_b^{(s)}$,
velocity sum $V_b^{(s)}=\sum_{i\in b}v_i^{(s)}$ and
squared-velocity sum $V_b^{2,(s)}=\sum_{i\in b}v_i^{(s)2}$ are
recorded. After averaging over realisations:
$\bar c_b = \langle c_b^{(s)}\rangle$,
$\bar V_b = \langle V_b^{(s)}\rangle$,
$\bar V_b^2 = \langle V_b^{2,(s)}\rangle$.
The hydrodynamic fields are
$\rho_b = \bar c_b / (M \Delta s)$,
$u_b = \bar V_b / \bar c_b$,
$T_b = \bar V_b^2/\bar c_b - u_b^2$.
Spatial and temporal derivatives use centred finite differences.

\begin{thebibliography}{99}

\bibitem{Tonks1936}
L.~Tonks,
Phys.~Rev.\ \textbf{50}, 955 (1936).

\bibitem{SalsburgZwanzig1953}
Z.~W.~Salsburg, R.~W.~Zwanzig, J.~G.~Kirkwood,
J.~Chem.~Phys.\ \textbf{21}, 1098 (1953).

\bibitem{LiebMattis1966}
B.~Sutherland, in
E.~H.~Lieb, D.~C.~Mattis (Eds.),
\emph{Mathematical Physics in One Dimension},
Academic Press (1966), Ch.~III.

\bibitem{Aizenman1975}
M.~Aizenman, S.~Goldstein, J.~L.~Lebowitz,
Commun.~Math.~Phys.\ \textbf{39}, 289 (1975).

\bibitem{Lebowitz1968}
J.~L.~Lebowitz, J.~K.~Percus, J.~Sykes,
Phys.~Rev.\ \textbf{171}, 224 (1968).

\bibitem{BDS1983}
C.~Boldrighini, R.~L.~Dobrushin, Yu.~M.~Sukhov,
J.~Stat.~Phys.\ \textbf{31}, 577 (1983).

\bibitem{Jepsen1965}
D.~W.~Jepsen,
J.~Math.~Phys.\ \textbf{6}, 405 (1965).

\bibitem{DoyonSpohn2017}
B.~Doyon, H.~Spohn,
J.~Stat.~Mech.\ 073210 (2017).

\bibitem{Doyon2017Lectures}
B.~Doyon,
SciPost Phys.\ Lect.\ Notes \textbf{18} (2020).

\bibitem{DenardisDoyon2018}
J.~De~Nardis, D.~Bernard, B.~Doyon,
Phys.~Rev.~Lett.\ \textbf{121}, 160603 (2018).

\bibitem{Bulchandani2018}
V.~B.~Bulchandani, R.~Vasseur, C.~Karrasch, J.~E.~Moore,
Phys.~Rev.~B \textbf{97}, 045407 (2018).

\bibitem{GHDperspective2025}
B.~Doyon, S.~Gopalakrishnan, F.~M{\o}ller,
J.~Schmiedmayer, R.~Vasseur,
Phys.~Rev.~X \textbf{15}, 010501 (2025).

\bibitem{HubnerDoyon2024}
F.~H\"ubner, B.~Doyon,
\emph{A new quadrature for the generalized hydrodynamics
equation and absence of shocks in the Lieb-Liniger model},
arXiv:2406.18322 (2024).

\bibitem{KhesinMisiolek2007}
B.~Khesin, G.~Misio{\l}ek,
Proc.~Steklov Inst.~Math.\ \textbf{259}, 73 (2007).

\bibitem{Lighthill1955}
M.~J.~Lighthill, G.~B.~Whitham,
Proc.~Roy.~Soc.~A \textbf{229}, 317 (1955).

\bibitem{Wei2000}
Q.-H.~Wei, C.~Bechinger, P.~Leiderer,
Science \textbf{287}, 625 (2000).

\bibitem{LinCui2015}
B.~Lin, B.~Cui, X.~Xu, et al.,
J.~Phys.\ Condens.~Matter \textbf{27}, 194116 (2015).

\bibitem{Schemmer2019}
M.~Schemmer, I.~Bouchoule, B.~Doyon, J.~Dubail,
Phys.~Rev.~Lett.\ \textbf{122}, 090601 (2019).

\bibitem{Malvania2021}
N.~Malvania, Y.~Zhang, Y.~Le, J.~Dubail, M.~Rigol, D.~S.~Weiss,
Science \textbf{373}, 1129 (2021).

\bibitem{Cataldini2022}
F.~Cataldini, F.~M{\o}ller, M.~Tajik et al.,
Phys.~Rev.~X \textbf{12}, 041032 (2022).

\bibitem{PaperI}
H.~Medel,
in preparation (2026).

\bibitem{PaperII}
H.~Medel,
in preparation (2026).

\end{thebibliography}

\end{document}
